\chapter{The Quantum Fourier Transform}

The Quantum Fourier Transform (QFT) is one of the most important subroutines in
quantum computing. It is the quantum analogue of the classical discrete Fourier
transform (DFT), and it sits at the heart of quantum phase estimation (QPE) and
of Shor's factoring algorithm. In this chapter we define the QFT, prove that it
is unitary, derive its product (tensor) representation, use that representation
to build an efficient circuit from Hadamard and controlled-phase gates, count
the gates to obtain the $O(n^2)$ bound, and discuss the inverse QFT together
with its principal applications.

Throughout we work with $n$ qubits, we write $N = 2^n$ for the dimension of the
state space $(\C^2)^{\otimes n}\cong\C^N$, and we abbreviate the primitive
$N$-th root of unity by
\[
  \omega \defeq e^{i 2\pi / N}, \qquad \omega^N = 1 .
\]
Computational basis states are labelled by integers $j\in\{0,1,\dots,N-1\}$, and
we freely identify the integer $j$ with its $n$-bit binary string
$j_1 j_2 \cdots j_n$, where $j_1$ is the most significant bit, so that
\[
  j = j_1 2^{n-1} + j_2 2^{n-2} + \cdots + j_n 2^{0}
    = \sum_{l=1}^{n} j_l\, 2^{\,n-l},
  \qquad \ket{j} = \ket{j_1}\otimes\cdots\otimes\ket{j_n}.
\]

\section{Definition and Motivation}

\begin{definition}[Quantum Fourier Transform]\label{l5:def:qft}
Let $N = 2^n$ and $\omega = e^{i 2\pi / N}$. The \emph{Quantum Fourier
Transform} is the linear operator $\QFT$ on $\C^N$ that acts on computational
basis states by
\[
  \QFT\ket{j} \;=\; \frac{1}{\sqrt N}\sum_{k=0}^{N-1} e^{i 2\pi j k / N}\ket{k}
             \;=\; \frac{1}{\sqrt N}\sum_{k=0}^{N-1} \omega^{jk}\,\ket{k},
  \qquad j = 0,1,\dots,N-1 .
\]
Equivalently, in the computational basis $\QFT$ is the $N\times N$ matrix with
entries
\[
  \bigl(\QFT\bigr)_{kj} \;=\; \frac{1}{\sqrt N}\,\omega^{jk}
  \qquad (0 \le j,k \le N-1),
\]
a (normalized) Vandermonde matrix in the powers of $\omega$.
\end{definition}

The action of $\QFT$ on a general state follows by linearity. If
$\ket{\psi} = \sum_{j=0}^{N-1} x_j \ket{j}$ then
\[
  \QFT\ket{\psi}
    = \sum_{k=0}^{N-1} y_k \ket{k},
  \qquad
  y_k = \frac{1}{\sqrt N}\sum_{j=0}^{N-1} \omega^{jk} x_j ,
\]
so the QFT transforms the vector of amplitudes $(x_0,\dots,x_{N-1})$ into
$(y_0,\dots,y_{N-1})$ exactly as the classical DFT would. The essential point is
that the input is a single computational basis state $\ket{j}$ while the output
is an \emph{equal-magnitude} superposition of all $N$ basis states; the integer
$j$ is encoded entirely in the relative phases $\omega^{jk}$.

\begin{example}[The all-zeros input]\label{l5:ex:zero}
For $j=0$ every phase is $\omega^{0}=1$, so
\[
  \QFT\ket{0} = \frac{1}{\sqrt N}\sum_{k=0}^{N-1}\ket{k}
              = H^{\otimes n}\ket{0}^{\otimes n},
\]
the uniform superposition. Thus on the input $\ket{0\cdots 0}$ the QFT is
nothing but a layer of Hadamard gates.
\end{example}

\begin{remark}[Same kernel as the classical DFT, but no free lunch]\label{l5:rem:dft}
The kernel $\tfrac{1}{\sqrt N}\,e^{i2\pi jk/N}$ appearing in
Definition~\ref{l5:def:qft} is \emph{identical} to that of the (symmetrically
normalized) classical DFT; the QFT simply applies this transform to the $2^n$
amplitudes of an $n$-qubit state \emph{in superposition}. It is tempting to
conclude that the QFT gives an exponential speedup for Fourier analysis, but
this is not so: after the transform we hold the coefficients $y_k$ only as
quantum amplitudes, and a measurement returns a \emph{single} basis state
$\ket{k}$ with probability $\abs{y_k}^2$. We cannot read out all $N$ amplitudes;
in fact estimating them to good precision would require exponentially many
repetitions. The power of the QFT is therefore not standalone but
\emph{structural}: embedded inside a larger algorithm (phase estimation,
period finding) it converts information hidden in the phases into information in
the measurement statistics, so that a small number of measurements suffices. We
also note the resource comparison developed below: the QFT circuit uses
$O(n^2) = O((\log N)^2)$ elementary gates, versus $O(N\log N)$ arithmetic
operations for the classical FFT and $O(N^2)$ for the naive DFT.
\end{remark}

\section{Unitarity}

Because the QFT is to be realized as a quantum circuit, it must be unitary. We
prove this directly from the definition. The engine of the proof is the
following elementary but crucial identity.

\begin{lemma}[Geometric-series / Kronecker-delta identity]\label{l5:lem:geom}
Let $N\ge 1$ and let $m\in\Z$. Then
\[
  \sum_{k=0}^{N-1} e^{i 2\pi m k / N}
  = \sum_{k=0}^{N-1}\bigl(e^{i2\pi m/N}\bigr)^{k}
  = N\,[\,m\equiv 0 \bmod N\,],
\]
where the bracket equals $1$ if $m\equiv 0 \pmod N$ and $0$ otherwise. In
particular, for $m = j - j'$ with $j,j' \in \{0,\dots,N-1\}$ we have
$m\equiv 0 \pmod N$ iff $j = j'$, so
\[
  \sum_{k=0}^{N-1} e^{i 2\pi (j-j') k / N} = N\,\delta_{j,j'} .
\]
\end{lemma}

\begin{proof}
Write $r = e^{i 2\pi m/N}$, so the sum is the geometric series
$\sum_{k=0}^{N-1} r^{k}$. If $m\equiv 0\pmod N$ then $r=1$ and every term is
$1$, giving the value $N$. Otherwise $r\ne 1$, and the closed form of the
geometric series applies:
\[
  \sum_{k=0}^{N-1} r^{k} = \frac{r^{N}-1}{r-1}
  = \frac{e^{i 2\pi m}-1}{r-1} = \frac{1-1}{r-1} = 0,
\]
since $r^{N} = e^{i 2\pi m} = 1$ for integer $m$ while the denominator is
nonzero. For the final statement, if $\abs{j-j'} < N$ then $j-j'\equiv 0
\pmod N$ forces $j-j'=0$.
\end{proof}

\begin{theorem}[The QFT is unitary]\label{l5:thm:unitary}
The operator of Definition~\ref{l5:def:qft} satisfies $\QFT^{\dagger}\QFT = \id$,
so $\QFT$ is unitary.
\end{theorem}

\begin{proof}
From Definition~\ref{l5:def:qft}, $\bigl(\QFT\bigr)_{kj} = N^{-1/2}\,\omega^{jk}$,
so the adjoint has entries
$\bigl(\QFT^{\dagger}\bigr)_{jk} = \overline{\bigl(\QFT\bigr)_{kj}}
= N^{-1/2}\,\omega^{-jk}$. We compute the matrix element
$\braket{j'|\QFT^{\dagger}\QFT|j}$, i.e.\ the $(j',j)$ entry of
$\QFT^{\dagger}\QFT$, by inserting a resolution of the identity over the
basis $\{\ket{k}\}$:
\[
  \braket{j'|\QFT^{\dagger}\QFT|j}
    = \sum_{k=0}^{N-1} \overline{(\QFT)_{kj'}}\,(\QFT)_{kj}
    = \sum_{k=0}^{N-1}
        \Bigl(\tfrac{1}{\sqrt N}\,\omega^{-j'k}\Bigr)
        \Bigl(\tfrac{1}{\sqrt N}\,\omega^{jk}\Bigr)
    = \frac{1}{N}\sum_{k=0}^{N-1}\omega^{(j-j')k}.
\]
Since $\omega^{(j-j')k} = e^{i2\pi (j-j')k/N}$, Lemma~\ref{l5:lem:geom} gives
$\sum_{k=0}^{N-1}\omega^{(j-j')k} = N\,\delta_{j,j'}$, and therefore
\[
  \braket{j'|\QFT^{\dagger}\QFT|j}
    = \frac{1}{N}\,N\,\delta_{j,j'} = \delta_{j,j'} .
\]
As $j,j'$ range over the basis this says $\QFT^{\dagger}\QFT = \id$. Equivalently,
the columns $\QFT\ket{j}$ are orthonormal. Hence $\QFT$ is unitary.
\end{proof}

Since $\QFT$ is unitary, $\QFT^{-1} = \QFT^{\dagger}$, and we record the inverse
explicitly.

\begin{definition}[Inverse QFT]\label{l5:def:iqft}
The \emph{inverse Quantum Fourier Transform} is $\QFT^{-1} = \QFT^{\dagger}$,
acting on basis states by
\[
  \QFT^{\dagger}\ket{j}
   = \frac{1}{\sqrt N}\sum_{k=0}^{N-1} e^{-i 2\pi j k / N}\ket{k}
   = \frac{1}{\sqrt N}\sum_{k=0}^{N-1} \omega^{-jk}\,\ket{k}.
\]
It differs from $\QFT$ only by conjugating every phase (replacing $\omega$ by
$\overline{\omega}=\omega^{-1}$). Consequently the inverse circuit is obtained
from the QFT circuit by reversing the order of the gates and negating every
phase angle.
\end{definition}

\section{The Product (Tensor) Representation}

To build a circuit we need the QFT output expressed qubit-by-qubit. The bridge
is the \emph{product form}, which shows that $\QFT\ket{j}$ is an unentangled
product state. First we fix notation for the binary fractions that appear in the
phases.

\begin{definition}[Binary fraction]\label{l5:def:binfrac}
For bits $j_l, j_{l+1}, \dots, j_n \in\{0,1\}$ the \emph{binary fraction}
$0.j_l j_{l+1}\cdots j_n$ is the real number
\[
  0.j_l j_{l+1}\cdots j_n
   \defeq \frac{j_l}{2} + \frac{j_{l+1}}{2^{2}} + \cdots + \frac{j_n}{2^{\,n-l+1}}
   = \sum_{s=l}^{n} \frac{j_s}{2^{\,s-l+1}} .
\]
In particular $0.j_n = j_n/2$ and $0.j_1 j_2\cdots j_n = j/2^{n} = j/N$.
\end{definition}

\begin{theorem}[Product representation]\label{l5:thm:product}
For every $j = j_1 j_2\cdots j_n$ (binary),
\[
  \QFT\ket{j}
   = \bigotimes_{l=1}^{n}
       \frac{1}{\sqrt 2}\Bigl(\ket{0}
        + e^{\,i 2\pi\, 0.j_l j_{l+1}\cdots j_n}\,\ket{1}\Bigr)
   = \frac{1}{2^{n/2}}
     \bigl(\ket{0}+e^{i2\pi\,0.j_1\cdots j_n}\ket{1}\bigr)\otimes\cdots\otimes
     \bigl(\ket{0}+e^{i2\pi\,0.j_n}\ket{1}\bigr) .
\]
In particular the output is a product state: the $n$ output qubits are
unentangled.
\end{theorem}

\begin{proof}
Start from Definition~\ref{l5:def:qft} and expand the summation index $k$ in binary,
$k = k_1 k_2\cdots k_n$ with $k = \sum_{l=1}^{n} k_l 2^{\,n-l}$. Then
\[
  \QFT\ket{j}
   = \frac{1}{\sqrt N}\sum_{k=0}^{N-1} e^{i2\pi jk/N}\ket{k}
   = \frac{1}{\sqrt N}
      \sum_{k_1=0}^{1}\cdots\sum_{k_n=0}^{1}
      e^{\,i 2\pi j \left(\sum_{l=1}^{n} k_l 2^{\,n-l}\right)/2^{n}}\,
      \ket{k_1}\cdots\ket{k_n}.
\]
The exponent factorizes over $l$ because the exponential of a sum is a product
of exponentials:
\[
  e^{\,i 2\pi j \left(\sum_{l} k_l 2^{\,n-l}\right)/2^{n}}
   = \prod_{l=1}^{n} e^{\,i 2\pi j k_l 2^{\,n-l}/2^{n}}
   = \prod_{l=1}^{n} e^{\,i 2\pi j k_l / 2^{\,l}} .
\]
Substituting and interchanging the (finite) sums with the product turns the
sum over all strings $k$ into a product of independent single-qubit sums:
\[
  \QFT\ket{j}
   = \frac{1}{\sqrt N}\bigotimes_{l=1}^{n}
       \Bigl(\sum_{k_l=0}^{1} e^{\,i 2\pi j k_l / 2^{\,l}}\ket{k_l}\Bigr)
   = \frac{1}{\sqrt N}\bigotimes_{l=1}^{n}
       \Bigl(\ket{0} + e^{\,i 2\pi j / 2^{\,l}}\ket{1}\Bigr).
\]
It remains to simplify the phase $e^{\,i 2\pi j/2^{\,l}}$. Write $j$ in binary and
divide by $2^{l}$:
\[
  \frac{j}{2^{l}}
   = \frac{1}{2^{l}}\sum_{s=1}^{n} j_s 2^{\,n-s}
   = \underbrace{\sum_{s\,:\,n-s\ge l} j_s 2^{\,n-s-l}}_{\text{integer}}
     \; + \; \sum_{s\,:\,n-s< l} j_s 2^{\,n-s-l}.
\]
The first group consists of nonnegative integer powers of two and contributes an
integer to $j/2^{l}$; since $e^{i2\pi(\text{integer})}=1$, it may be discarded.
The second group runs over $s = n-l+1, n-l+2, \dots, n$ and gives, after
relabelling,
\[
  \sum_{s=n-l+1}^{n} j_s 2^{\,n-s-l}
   = \frac{j_{\,n-l+1}}{2} + \frac{j_{\,n-l+2}}{2^{2}} + \cdots + \frac{j_n}{2^{\,l}}
   = 0.j_{\,n-l+1}\,j_{\,n-l+2}\cdots j_n .
\]
Hence $e^{\,i 2\pi j/2^{\,l}} = e^{\,i 2\pi\, 0.j_{\,n-l+1}\cdots j_n}$. Finally,
$1/\sqrt N = 2^{-n/2} = \prod_{l=1}^n 2^{-1/2}$ distributes one factor
$1/\sqrt2$ into each tensor factor, giving
\[
  \QFT\ket{j}
   = \bigotimes_{l=1}^{n}\frac{1}{\sqrt2}
     \Bigl(\ket{0}+e^{\,i2\pi\,0.j_{\,n-l+1}\cdots j_n}\ket{1}\Bigr).
\]
Reindexing the product by $l' = n-l+1$ (so that the factor with the shortest
binary fraction $0.j_n$ moves to one end and the longest $0.j_1\cdots j_n$ to the
other) puts the factors in the standard order stated in the theorem, in which
the first output qubit carries the phase $0.j_1\cdots j_n$ and the last carries
$0.j_n$. The output is a tensor product of single-qubit states, hence
unentangled.
\end{proof}

\section{The Circuit}

The product representation immediately suggests how to synthesize the QFT one
qubit at a time. Two gate families suffice: the Hadamard gate $H$ and the
controlled-phase gates $R_k$.

\begin{definition}[Controlled-phase gate $R_k$]\label{l5:def:Rk}
For an integer $k\ge 1$ let $R_k$ be the single-qubit phase gate
\[
  R_k = \begin{pmatrix} 1 & 0 \\ 0 & e^{i 2\pi/2^{k}} \end{pmatrix}
      = \operatorname{diag}\!\bigl(1,\; e^{i 2\pi/2^{k}}\bigr).
\]
Thus $R_1 = \operatorname{diag}(1,-1) = Z$, $R_2 = \operatorname{diag}(1,i) = S$
(the $\pi/2$ phase gate), $R_3 = \operatorname{diag}(1, e^{i\pi/4}) = T$, and so
on. In the circuit we use the \emph{controlled} version, in which $R_k$ is
applied to a target qubit conditioned on a control qubit being $\ket{1}$.
\end{definition}

The construction rests on the single-qubit identity, valid for a bit
$j_m\in\{0,1\}$,
\[
  H\ket{j_m}
   = \frac{1}{\sqrt2}\bigl(\ket{0} + e^{i\pi j_m}\ket{1}\bigr)
   = \frac{1}{\sqrt2}\bigl(\ket{0} + e^{i 2\pi\, 0.j_m}\ket{1}\bigr),
\]
because $e^{i\pi j_m} = (-1)^{j_m}$ and $0.j_m = j_m/2$. A Hadamard therefore
produces the leading term $0.j_m$ of the required phase; the lower-order bits
are then folded in by controlled $R_k$ gates, using the split
\[
  e^{i2\pi\,0.j_m j_{m+1}\cdots j_n}
   = e^{i2\pi\,0.j_m}\cdot e^{i2\pi\,0.0 j_{m+1}}\cdots e^{i2\pi\,0.0\cdots 0 j_n},
\]
where the factor $e^{i2\pi\,0.\overbrace{0\cdots0}^{k-1}\,j_{m+k-1}}
= e^{i2\pi j_{m+k-1}/2^{k}}$ is exactly the phase $R_k$ applies when its control
qubit holds the bit $j_{m+k-1}$.

The full circuit processes the qubits in order $1,2,\dots,n$. On qubit $m$ we
apply $H$, and then for each $k = 2,3,\dots, n-m+1$ a controlled-$R_k$ with the
target being qubit $m$ and the control being qubit $m+k-1$. After all qubits are
processed, the resulting single-qubit phases appear on the qubits in reversed
order relative to Theorem~\ref{l5:thm:product}; a final layer of SWAP gates
(\emph{bit reversal}) restores the correct qubit order.

\begin{example}[The $n=2$ QFT circuit, worked in full]\label{l5:ex:n2}
Take $n=2$, $N=4$, so the input is $\ket{j} = \ket{j_1 j_2}$ (with $j_1$ the
high bit). We track the state through three steps and a final swap.

\emph{Step 1 (Hadamard on qubit $1$).} Applying $H\otimes\id$,
\[
  (H\otimes \id)\ket{j_1 j_2}
   = \frac{1}{\sqrt2}\bigl(\ket{0}+e^{i2\pi\,0.j_1}\ket{1}\bigr)\otimes\ket{j_2}.
\]
The first qubit now carries the partial phase $0.j_1$; it still needs the
contribution of $j_2$.

\emph{Step 2 (controlled $R_2$, control $=$ qubit $2$, target $=$ qubit $1$).}
With $R_2 = \operatorname{diag}(1, e^{i\pi/2})$, the controlled gate multiplies
the $\ket{1}$ component of qubit~$1$ by $e^{i2\pi j_2/2^{2}} = e^{i2\pi\,0.0 j_2}$,
upgrading the phase $0.j_1$ to the full $0.j_1 j_2$:
\[
   \frac{1}{\sqrt2}\bigl(\ket{0}+e^{i2\pi\,0.j_1 j_2}\ket{1}\bigr)\otimes\ket{j_2}.
\]

\emph{Step 3 (Hadamard on qubit $2$).} Applying $\id\otimes H$ turns
$\ket{j_2}$ into $\tfrac{1}{\sqrt2}(\ket{0}+e^{i2\pi\,0.j_2}\ket{1})$, giving
\[
   \frac{1}{2}
     \bigl(\ket{0}+e^{i2\pi\,0.j_1 j_2}\ket{1}\bigr)\otimes
     \bigl(\ket{0}+e^{i2\pi\,0.j_2}\ket{1}\bigr).
\]

\emph{Final SWAP (bit reversal).} Comparing with Theorem~\ref{l5:thm:product}, the
QFT requires the first output qubit to carry $0.j_1 j_2$ and the second to carry
$0.j_2$. The circuit above has produced the phase $0.j_1 j_2$ on the line that
started as qubit~$1$ and $0.j_2$ on the line that started as qubit~$2$; in the
general $n$-qubit construction these phases emerge in reversed order, so a SWAP
layer is required. For $n=2$ a single SWAP exchanges the two lines to yield the
standard ordering
\[
   \QFT\ket{j}
    = \frac{1}{2}
      \bigl(\ket{0}+e^{i2\pi\,0.j_1 j_2}\ket{1}\bigr)\otimes
      \bigl(\ket{0}+e^{i2\pi\,0.j_2}\ket{1}\bigr),
\]
which is exactly Theorem~\ref{l5:thm:product} for $n=2$. In summary the two-qubit
QFT is: $H$ on qubit $1$; controlled-$R_2$ from qubit $2$ onto qubit $1$; $H$ on
qubit $2$; then SWAP.

\emph{A concrete input.} Take $j = 3 = (11)_2$, so $j_1=j_2=1$. Here
$0.j_2 = 0.1 = 1/2$ and $0.j_1 j_2 = 0.11 = 3/4$, hence
$e^{i2\pi\,0.j_1 j_2} = e^{i 3\pi/2} = -i$ and $e^{i2\pi\,0.j_2} = e^{i\pi} = -1$.
Therefore
\[
  \QFT\ket{3}
   = \frac{1}{2}\bigl(\ket{0}- i\ket{1}\bigr)\otimes\bigl(\ket{0}-\ket{1}\bigr)
   = \frac{1}{2}\bigl(\ket{0}-\ket{1}- i\ket{2}+ i\ket{3}\bigr).
\]
This agrees with Definition~\ref{l5:def:qft}: for $N=4$, $\omega = e^{i\pi/2}=i$,
so
\[
  \QFT\ket{3} = \tfrac12\sum_{k=0}^{3} i^{3k}\ket{k}
   = \tfrac12\bigl(\ket{0}+i^{3}\ket{1}+i^{6}\ket{2}+i^{9}\ket{3}\bigr)
   = \tfrac12\bigl(\ket{0}-i\ket{1}-\ket{2}+i\ket{3}\bigr),
\]
which matches the tensor-product result once the bit-reversal SWAP has aligned
the qubit order with the integer labels $\ket{0},\ket{1},\ket{2},\ket{3}$.
\end{example}

\section{Gate Count}

\begin{theorem}[Gate complexity of the QFT]\label{l5:thm:count}
The circuit described above realizes the QFT on $n$ qubits using exactly $n$
Hadamard gates, $\binom{n}{2} = \tfrac{n(n-1)}{2}$ controlled-phase gates, and at
most $\lfloor n/2\rfloor$ SWAP gates. The total is
\[
  n + \frac{n(n-1)}{2} + \Bigl\lfloor\frac n2\Bigr\rfloor = O(n^{2})
\]
elementary gates. Since $n = \log_2 N$, this is $O\bigl((\log N)^{2}\bigr)$,
i.e.\ polylogarithmic in the dimension $N$.
\end{theorem}

\begin{proof}
Process qubits in order $m = 1,2,\dots,n$. On qubit $m$ we apply exactly one
Hadamard, contributing $n$ Hadamards in total. After the Hadamard on qubit $m$,
we must fold in the bits $j_{m+1},\dots,j_n$, one controlled-phase gate each,
namely controlled-$R_2,\dots,$ controlled-$R_{\,n-m+1}$; that is $n-m$
controlled-phase gates on qubit $m$. Summing over $m$,
\[
  \sum_{m=1}^{n} (n-m) = (n-1)+(n-2)+\cdots+1+0
   = \frac{n(n-1)}{2} = \binom{n}{2},
\]
which is $O(n^2)$. Finally, bit reversal exchanges qubit $m$ with qubit $n-m+1$
for $m=1,\dots,\lfloor n/2\rfloor$, using $\lfloor n/2\rfloor = O(n)$ SWAP gates
(each SWAP is three CNOTs, still $O(n)$ two-qubit gates). Adding the three
contributions gives $n + \tfrac{n(n-1)}{2} + \lfloor n/2\rfloor = O(n^2)$.
\end{proof}

\begin{remark}[Approximate QFT and comparison with classical transforms]\label{l5:rem:approx}
The controlled-$R_k$ gates apply a phase $e^{i2\pi/2^{k}}$ that becomes
exponentially small in $k$: for large $k$ these gates are nearly the identity
and are the hardest to implement accurately on hardware. Dropping all gates with
$k > k_{\max} = O(\log(n/\varepsilon))$ yields the \emph{approximate QFT}, which
is within error $\varepsilon$ of the exact QFT in operator norm while using only
$O(n\log(n/\varepsilon))$ gates. Either way, the QFT's $O(n^2)$ (or
$O(n\log n)$) gate count is dramatically smaller than the $O(N\log N)$
arithmetic operations of the classical Fast Fourier Transform and the $O(N^2)$
of the naive DFT --- recalling that $N = 2^n$. As stressed in
Remark~\ref{l5:rem:dft}, this gate-count advantage is not by itself an end-to-end
speedup, because a single measurement extracts only $O(n)$ classical bits from
the transformed state.
\end{remark}

\section{Applications}

\paragraph{Phase estimation (QPE).}
Suppose a unitary $U$ has eigenvector $\ket{u}$ with eigenvalue
$e^{i 2\pi\varphi}$, where $\varphi\in[0,1)$ has the exact $n$-bit binary
expansion $\varphi = 0.\varphi_1\varphi_2\cdots\varphi_n$. Using controlled
powers $U^{2^{t}}$ acting on a Hadamard-initialized $n$-qubit register, one
prepares the state
\[
  \frac{1}{2^{n/2}}\bigotimes_{l=1}^{n}
    \bigl(\ket{0}+e^{i2\pi\,0.\varphi_l\cdots\varphi_n}\ket{1}\bigr)
  = \QFT\ket{\varphi_1\cdots\varphi_n},
\]
which by Theorem~\ref{l5:thm:product} is precisely $\QFT$ applied to the basis
state encoding $\varphi$. Applying the \emph{inverse} QFT and measuring in the
computational basis therefore returns the bits $\varphi_1\cdots\varphi_n$
exactly (and, when $\varphi$ is not an exact $n$-bit fraction, its best $n$-bit
approximation with high probability). This is the subject of the next lecture.

\paragraph{Shor's period finding.}
For a periodic function $f$ with unknown period $r$, preparing a superposition
over the domain and applying $\Uf$ produces a state that is periodic in the
computational register. Applying the QFT concentrates the amplitude on basis
states $k$ near integer multiples of $N/r$: the hidden period is converted into
sharp \emph{measurable} peaks whose positions, via continued fractions, reveal
$r$. This period-finding routine is the quantum core of Shor's factoring and
discrete-logarithm algorithms. In both applications the QFT (or its inverse)
acts as a general-purpose engine that reads phase-encoded information out into
measurement statistics.

\section{Problems}

\begin{exercise}\label{l5:ex:p1}
Compute $\QFT\ket{1}$ for $n=2$ ($N=4$) directly from Definition~\ref{l5:def:qft},
and verify the result using the product representation
(Theorem~\ref{l5:thm:product}).
\end{exercise}

\begin{exercise}\label{l5:ex:p2}
Prove that $\QFT\,\QFT^{\dagger} = \id$ as well, i.e.\ that $\QFT$ is unitary from
both sides, by evaluating $\braket{k'|\QFT\,\QFT^{\dagger}|k}$ with the identity
of Lemma~\ref{l5:lem:geom}.
\end{exercise}

\begin{exercise}\label{l5:ex:p3}
Show that for the single-qubit case $n=1$ the QFT equals the Hadamard gate,
$\QFT = H$. (Here $N=2$ and $\omega = e^{i\pi} = -1$.)
\end{exercise}

\begin{exercise}\label{l5:ex:p4}
Using the geometric-series identity, show that $\QFT$ maps the uniform
superposition $\tfrac{1}{\sqrt N}\sum_{j=0}^{N-1}\ket{j}$ to the basis state
$\ket{0}$. Interpret this as the statement that a constant signal has only a
zero-frequency component.
\end{exercise}

\begin{exercise}\label{l5:ex:p5}
Write down, gate by gate, the $n=3$ QFT circuit (Hadamards, controlled-$R_k$
gates, and SWAPs), and count the total number of gates. Confirm that it matches
Theorem~\ref{l5:thm:count}.
\end{exercise}

\begin{exercise}\label{l5:ex:p6}
For $N=2^n$ let $\ket{\psi_a} = \QFT\ket{a}$. Show that the states
$\{\ket{\psi_a}\}_{a=0}^{N-1}$ form an orthonormal basis, and compute the inner
product $\braket{\psi_a|\psi_b}$ for general $a,b$.
\end{exercise}

\begin{exercise}\label{l5:ex:p7}
Consider the ``shift by $1$'' operator $V$ defined on the computational basis by
$V\ket{j} = \ket{(j+1)\bmod N}$. Show that the QFT diagonalizes $V$: prove that
each $\ket{\psi_k} \defeq \QFT\ket{k}$ satisfies
$V\ket{\psi_k} = e^{-i2\pi k/N}\ket{\psi_k}$, so that the QFT basis states are
eigenvectors of $V$.
\end{exercise}

\begin{solution}
\textit{Exercise~\ref{l5:ex:p1}.} For $N=4$, $\omega=e^{i2\pi/4}=i$. By
Definition~\ref{l5:def:qft},
\[
  \QFT\ket{1} = \tfrac12\sum_{k=0}^{3} i^{k}\ket{k}
   = \tfrac12\bigl(\ket{0}+i\ket{1}-\ket{2}-i\ket{3}\bigr).
\]
For the product form, $j=1=(01)_2$ means $j_1=0,\ j_2=1$, so $0.j_2 = 0.1 =
\tfrac12$ and $0.j_1 j_2 = 0.01 = \tfrac14$, giving phases
$e^{i2\pi/4}=i$ and $e^{i2\pi/2}=-1$. Hence
$\QFT\ket{1} = \tfrac12(\ket{0}+i\ket{1})\otimes(\ket{0}-\ket{1})
= \tfrac12(\ket{0}-\ket{1}+i\ket{2}-i\ket{3})$, which equals the direct result
after the bit-reversal reordering that identifies the tensor-product qubit order
with the integer labels.

\medskip
\textit{Exercise~\ref{l5:ex:p2}.} With $(\QFT)_{kj}=N^{-1/2}\omega^{jk}$ and
$(\QFT^{\dagger})_{jk}=N^{-1/2}\omega^{-jk}$,
\[
  \braket{k'|\QFT\,\QFT^{\dagger}|k}
   = \sum_{j=0}^{N-1}(\QFT)_{k'j}\,(\QFT^{\dagger})_{jk}
   = \frac{1}{N}\sum_{j=0}^{N-1}\omega^{k'j}\omega^{-jk}
   = \frac{1}{N}\sum_{j=0}^{N-1}\omega^{j(k'-k)}
   = \delta_{k,k'},
\]
using Lemma~\ref{l5:lem:geom}. Thus $\QFT\,\QFT^{\dagger}=\id$; combined with
Theorem~\ref{l5:thm:unitary} this confirms $\QFT$ is unitary.

\medskip
\textit{Exercise~\ref{l5:ex:p3}.} For $n=1$, $N=2$ and $\omega=e^{i\pi}=-1$, so
$(\QFT)_{kj}=\tfrac{1}{\sqrt2}(-1)^{jk}$, i.e.\ the matrix
$\tfrac{1}{\sqrt2}\left(\begin{smallmatrix}1&1\\1&-1\end{smallmatrix}\right)=H$.
Concretely $\QFT\ket{0}=\tfrac{1}{\sqrt2}(\ket0+\ket1)$ and
$\QFT\ket{1}=\tfrac{1}{\sqrt2}(\ket0-\ket1)$, matching $H$.

\medskip
\textit{Exercise~\ref{l5:ex:p4}.} By linearity,
\[
  \QFT\Bigl(\tfrac{1}{\sqrt N}\sum_{j} \ket{j}\Bigr)
   = \frac{1}{N}\sum_{k}\Bigl(\sum_{j}\omega^{jk}\Bigr)\ket{k}
   = \frac{1}{N}\sum_{k} N\delta_{k,0}\ket{k} = \ket{0},
\]
where $\sum_j \omega^{jk} = N\delta_{k,0}$ by Lemma~\ref{l5:lem:geom}. A constant
(``DC'') signal has all its weight at zero frequency, $k=0$.

\medskip
\textit{Exercise~\ref{l5:ex:p5}.} Processing qubits $1,2,3$:
qubit~$1$: $H$, then controlled-$R_2$ (control $2$), controlled-$R_3$ (control
$3$); qubit~$2$: $H$, then controlled-$R_2$ (control $3$); qubit~$3$: $H$;
finally SWAP qubits $1$ and $3$. Counting: $3$ Hadamards, $2+1+0 = 3 =
\binom{3}{2}$ controlled-phase gates, and $1$ SWAP ($\lfloor 3/2\rfloor = 1$),
for $3+3+1=7$ gates, in agreement with Theorem~\ref{l5:thm:count}
($n+\tfrac{n(n-1)}2+\lfloor n/2\rfloor = 3+3+1$).

\medskip
\textit{Exercise~\ref{l5:ex:p6}.} Since $\QFT$ is unitary
(Theorem~\ref{l5:thm:unitary}) it maps the orthonormal basis $\{\ket{a}\}$ to the
orthonormal set $\{\ket{\psi_a}\}$, so the latter is a basis. Explicitly,
\[
  \braket{\psi_a|\psi_b}
   = \braket{a|\QFT^{\dagger}\QFT|b} = \braket{a|b} = \delta_{a,b},
\]
by Theorem~\ref{l5:thm:unitary}. (One may also verify this directly with
Lemma~\ref{l5:lem:geom}: $\braket{\psi_a|\psi_b} = \tfrac1N\sum_k
\omega^{k(b-a)} = \delta_{a,b}$.)

\medskip
\textit{Exercise~\ref{l5:ex:p7}.} Apply $V$ to a QFT basis state:
\[
  V\ket{\psi_k}
   = V\,\frac{1}{\sqrt N}\sum_{j}\omega^{jk}\ket{j}
   = \frac{1}{\sqrt N}\sum_{j}\omega^{jk}\ket{(j+1)\bmod N}.
\]
Reindex with $j' = (j+1)\bmod N$, i.e.\ $j = (j'-1)\bmod N$; because
$\omega^{N}=1$ the phase $\omega^{jk} = \omega^{(j'-1)k} =
\omega^{-k}\omega^{j'k}$ is well defined modulo $N$. Hence
\[
  V\ket{\psi_k} = \omega^{-k}\,\frac{1}{\sqrt N}\sum_{j'}\omega^{j'k}\ket{j'}
   = \omega^{-k}\ket{\psi_k} = e^{-i2\pi k/N}\ket{\psi_k}.
\]
Thus $\ket{\psi_k}=\QFT\ket{k}$ is an eigenvector of $V$ with eigenvalue
$e^{-i2\pi k/N}$, and equivalently $\QFT^{\dagger} V\,\QFT\ket{k} =
e^{-i2\pi k/N}\ket{k}$, so $\QFT^{\dagger} V\QFT$ is diagonal: the QFT
diagonalizes the cyclic shift.
\end{solution}
